% ==============================================================================
% MÓDULO: REFERENCIAS BIBLIOGRÁFICAS (Normas Vancouver NLM / ICMJE)
% Archivo: subcapitulos/sec_06_referencias.tex
% Orden estricto de primera aparición en Introducción, Capítulo I y Capítulo II
% Total de referencias activas: 20
% ==============================================================================
\section*{REFERENCIAS BIBLIOGRÁFICAS}
\label{sec:referencias}
\addcontentsline{toc}{section}{REFERENCIAS BIBLIOGRÁFICAS}

\begin{thebibliography}{99}
\itemsep 4pt

\bibitem{who2018}
World Health Organization, OECD, The World Bank. Delivering quality health services: a global imperative for universal health coverage. Geneva: World Health Organization; 2018.

\bibitem{minsacalidad2021}
Ministerio de Salud del Perú. Política Nacional de Calidad en Salud: Documento Técnico. Lima: MINSA; 2021.

\bibitem{minsainv2025}
Ministerio de Salud del Perú. Líneas de Investigación en Salud al 2030: Resolución Ministerial N° 424-2025/MINSA. Lima: MINSA; 2025.

\bibitem{diresa2025}
Dirección Regional de Salud de Ayacucho. Prioridades Regionales de Investigación en Salud de Ayacucho 2025--2030. Ayacucho: DIRESA Ayacucho; 2025.

\bibitem{sampieri2018}
Hernández-Sampieri R, Mendoza CP. Metodología de la investigación: las rutas cuantitativa, cualitativa y mixta. Ciudad de México: McGraw-Hill; 2018.

\bibitem{creswell2013}
Creswell JW. Qualitative inquiry and research design: Choosing among five approaches. 3rd ed. Thousand Oaks (CA): SAGE Publications; 2013.

\bibitem{husserl1949}
Husserl E. Ideas relativas a una fenomenología pura y una filosofía fenomenológica. Gaos J, traductor. México D.F.: Fondo de Cultura Económica; 1949.

\bibitem{vanmanen1990}
van Manen M. Researching lived experience: Human science for an action sensitive pedagogy. Albany (NY): State University of New York Press; 1990.

\bibitem{donabedian1980}
Donabedian A. Explorations in Quality Assessment and Monitoring. Vol. 1: The Definition of Quality and Approaches to its Assessment. Ann Arbor (MI): Health Administration Press; 1980.

\bibitem{donabedian2005}
Donabedian A. Evaluating the quality of medical care. 1966. Milbank Q. 2005;83(4):691-729.

\bibitem{penchansky1981}
Penchansky R, Thomas JW. The concept of access: definition and relationship to consumer satisfaction. Med Care. 1981;19(2):127-40.

\bibitem{watson2008}
Watson J. Nursing: The philosophy and science of caring. 2nd ed. Boulder: University Press of Colorado; 2008.

\bibitem{gubalincoln1985}
Lincoln YS, Guba EG. Naturalistic Inquiry. Beverly Hills (CA): SAGE Publications; 1985.

\bibitem{okuda2005}
Okuda Benavides M, Gómez-Restrepo C. Métodos en investigación cualitativa: triangulación. Rev Colomb Psiquiatr. 2005;34(1):118-24.

\bibitem{bunge2014}
Bunge M. La ciencia, su método y su filosofía. Buenos Aires: Editorial Sudamericana; 2014.

\bibitem{helsinki2013}
Asociación Médica Mundial. Declaración de Helsinki de la AMM: Principios éticos para las investigaciones médicas en seres humanos. Ferney-Voltaire: WMA; 2013.

\bibitem{cioms2016}
Consejo de Organizaciones Internacionales de las Ciencias Médicas (CIOMS). Pautas éticas internacionales para la investigación relacionada con la salud con seres humanos. Ginebra: CIOMS; 2016.

\bibitem{aguirre2020}
Aguirre-Andrade M, Tenorio-Acosta I, Rivas-Díaz L, Valenzuela-Oré F, Sánchez-Simbrón M, Quino-Huamaní F. Promoción de salud y nivel de salubridad de las familias, en comunidades del Distrito de Cangallo, Ayacucho 2020. Rev Investig Salud. 2020;14(2):45-58.

\bibitem{parasuraman1988}
Parasuraman A, Zeithaml VA, Berry LL. SERVQUAL: A multiple-item scale for measuring consumer perceptions of service quality. J Retailing. 1988;64(1):12-40.

\bibitem{minsaservqual2011}
Ministerio de Salud del Perú. Guía técnica para la evaluación de la satisfacción del usuario externo en los establecimientos de salud y servicios médicos de apoyo: Resolución Ministerial N° 527-2011/MINSA. Lima: MINSA; 2011.

\end{thebibliography}

\newpage
